\documentclass[11pt]{article}

\usepackage[margin=1in]{geometry}
\usepackage{amsmath,amssymb}
\usepackage{enumitem}
\usepackage{hyperref}
\hypersetup{hidelinks}

\title{Scaling Book Exercises: Section 4}
\author{Lucas Yan}
\date{}

\newcommand{\AI}{\operatorname{AI}}

\begin{document}
\maketitle

\noindent
\textbf{Chapter:} All the Transformer Math You Need to Know\\
\textbf{Source scan:} \texttt{../handwritten/section\_04\_handwritten.pdf}

\paragraph{Notation.}
$B$ is batch size, $T$ is query length, $S$ is KV length, $D=d_{\mathrm{model}}$,
$F=d_{\mathrm{ff}}$, $L$ is layer count, $V$ is vocabulary size, $N$ is the
number of query heads, $K$ is the number of KV heads, $H$ is head dimension,
and $G=N/K$ is the number of query heads per KV head. Thus $NH=D$.

\section*{Exercise 1}

For $D=4096$, $F=4D=16384$, $V=32000$, and $L=64$, the parameter count is
\[
P=L(3DF+4DNH+D)+2DV.
\]
Using $NH=D$,
\begin{align*}
P
&=64\bigl(3(4096)(16384)+4(4096)^2+4096\bigr)
  +2(4096)(32000)\\
&=16{,}001{,}335{,}296
 \approx 16.0\text{B parameters}.
\end{align*}

Ignoring the small normalization and embedding terms, the attention fraction is
\[
\frac{4D^2}{4D^2+3D(4D)}=\frac{4}{16}=\frac14.
\]

With int8 keys and values, the KV cache per token is
\[
2LNH=2LD=2(64)(4096)=524{,}288\text{ bytes}=512\text{ KiB}.
\]

\section*{Exercise 2}

For $A[B_X,D_Y] *_D W[D_Y,F]$ on a mesh
$\{X:4,Y:8,Z:4\}$, the mathematical operation has
\[
2BDF
\]
FLOPs. Since the computation is replicated across the unsharded $Z$ axis, the
hardware performs
\[
2BDFZ
\]
FLOPs in total. Each device handles a $1/(XY)$ shard of the mathematical work,
so the per-device work is
\[
\frac{2BDF}{XY}=\frac{2BDF}{32}=\frac{BDF}{16}.
\]

\section*{Exercise 3}

In
\[
A[I,J,K,L] * B[I,J,M,N,O] \longrightarrow C[K,L,M,N,O],
\]
$I$ and $J$ are contracting dimensions. Every distinct dimension contributes
once, and multiply-add contributes a factor of two. Therefore,
\[
\boxed{2IJKLMNO\text{ FLOPs}}.
\]

\section*{Exercise 4}

Ignoring the Q/K/V/O projections and softmax's lower-order cost, grouped-query
self-attention reads Q and the KV cache, performs the $QK^\top$ and $AV$
matmuls, and writes the output. In bf16,
\begin{align*}
\text{FLOPs} &= 4BTSNH=4BTSKGH,\\
\text{bytes} &=4BTNH+4BSKH=4BHK(TG+S).
\end{align*}
Hence
\[
\boxed{\AI(T,S,G)=\frac{TSG}{TG+S}}.
\]

During prefill, $S=T$, so
\[
\AI_{\mathrm{prefill}}=\frac{TG}{G+1}.
\]
On TPU v5e, whose hardware arithmetic intensity is approximately
$240$ FLOPs/byte, attention becomes compute-bound when
\[
T\gtrsim 240\frac{G+1}{G}.
\]
For MHA, $G=1$, giving $T\gtrsim480$; for large $G$, the threshold approaches
$240$.

During single-token generation, $T=1$, and
\[
\AI_{\mathrm{generation}}=\frac{SG}{G+S}\longrightarrow G
\quad(S\to\infty).
\]
Thus generation attention is normally bandwidth-bound because $G\ll240$.

For prefill with $S=T$, the attention-to-FFW FLOP ratio is
\[
\frac{12BT^2NH}{18BTDF}
=\frac{2T}{3F}
=\frac{T}{6D}\qquad(F=4D).
\]
After accounting for the bandwidth roofline, attention's effective cost is
nearly flat while bandwidth-bound, then grows linearly after the crossover.

\section*{Exercise 5}

Training FLOPs for the Q/K/V/O projections and self-attention are respectively
\[
24BTDNH
\quad\text{and}\quad
12BT^2NH.
\]
Equating them gives
\[
24BTDNH=12BT^2NH
\Longrightarrow T=2D.
\]
For $D=4096$,
\[
\boxed{T=8192}.
\]

\section*{Exercise 6}

If only the seven principal matmul outputs are saved, the backward pass must
recompute
\[
QK^\top
\quad\text{and}\quad
\operatorname{softmax}(QK^\top)V
\]
to obtain the output-projection gradient. Each matmul costs $2BT^2NH$, so the
additional matmul work is
\[
\boxed{4BT^2NH\text{ FLOPs}}.
\]
There are also lower-order elementwise recomputations of order $O(BTD)$ and
$O(BTF)$.

\section*{Exercise 7}

The H800's FP8 peak without structured sparsity is approximately half of the
advertised sparse value:
\[
\frac{3026}{2}=1513\text{ TFLOPs/s}=1.513\times10^{15}\text{ FLOPs/s}.
\]
The available hardware FLOPs over $2.79$ million H800-hours are
\[
(2.79\times10^6)(1.513\times10^{15})(3600)
\approx1.52\times10^{25}.
\]
Using the $6NT$ training approximation with $37$B activated parameters and
$14.8$T tokens,
\[
6(37\times10^9)(14.8\times10^{12})
\approx3.29\times10^{24}\text{ FLOPs}.
\]
Therefore the model FLOPs utilization is
\[
\boxed{\frac{3.29\times10^{24}}{1.52\times10^{25}}\approx21.7\%}.
\]

\section*{Exercise 8}

For one int8 expert matrix, loading all experts costs $EDF$ bytes. If every
token selects $k$ experts, the computation costs $2kBDF$ FLOPs. Thus
\[
\AI=\frac{2kBDF}{EDF}=\frac{2kB}{E}.
\]
To exceed TPU v5e's approximate $240$ FLOPs/byte roofline,
\[
\frac{2kB}{E}>240
\Longrightarrow
\boxed{B>\frac{120E}{k}}.
\]
For DeepSeek with $E=256$ routed experts and $k=8$,
\[
\boxed{B>\frac{120(256)}{8}=3840\text{ tokens}}.
\]

\end{document}
